\documentclass[12pt]{article}

\usepackage{amsmath,amssymb,amsthm}
\usepackage{bm}
\usepackage{hyperref}
\usepackage{mathtools}

% Middle-Way operator: \mweq
\DeclareMathOperator{\mw}{mw}
\newcommand{\mweq}{\mathrel{\overset{\mathrm{mw}}{=}}}


\title{An Extended Resolution of the Classical P=NP Problem via the Middle-Way Equivalence Operator $\mweq$}
\author{Masaomi Chiba}
\date{2025-12-10}

\newtheorem{definition}{Definition}
\newtheorem{theorem}{Theorem}
\newtheorem{lemma}{Lemma}
\newtheorem{proposition}{Proposition}
\newtheorem{remark}{Remark}


\begin{document}

\maketitle

\begin{abstract}
The classical formulation of the P=NP problem presupposes a \emph{static notion of equality},
in which the equality symbol ``$=$'' denotes strict identity between complexity classes. 
However, empirical computational practice---from SAT solvers and heuristic algorithms 
to machine learning and large language models---reveals a different picture: 
NP-complete problems often behave as if they are ``practically tractable'' despite their 
worst-case hardness.

This paper introduces a new operator, the \emph{Middle-Way Equivalence} $\mweq$, 
which expresses \emph{dynamic, procedural, and convergence-based equivalence} 
between complexity classes. 
We show that under this operator, the relationship between P and NP can be expressed as
\[
\mathrm{P} \mweq \mathrm{NP},
\]
a formulation that captures the empirical convergence and practical equivalence 
between these classes without asserting classical equality.

This yields an \emph{extended resolution} of the P=NP problem, not by deciding 
its truth value in the classical sense, but by generalizing the notion of equivalence 
to account for real-world computational phenomena.
\end{abstract}

\section{Introduction}

The P=NP problem traditionally asks whether
\[
\mathrm{P} = \mathrm{NP},
\]
where ``$=$'' denotes strict set-theoretic equality. 
This classical framing assumes:

\begin{itemize}
    \item fixed computational models,
    \item static notions of identity,
    \item observer-independence of computational truth,
    \item and worst-case performance as the essential criterion.
\end{itemize}

However, modern computational systems---including heuristic search,
SAT solvers, probabilistic algorithms, evolutionary algorithms,
reinforcement learning, and large language models---exhibit behavior
that contradicts the naive expectation that NP-complete problems are ``intractable.''

In practice:

\begin{itemize}
    \item NP-complete problems are solved efficiently in most real-world cases;
    \item heuristic and approximate algorithms achieve near-optimal solutions;
    \item machine learning models infer NP solutions statistically;
    \item adaptive computation reduces effective complexity.
\end{itemize}

These observations suggest a need to reformulate the equivalence relation 
used in the P=NP problem.

\section{Motivation for the Middle-Way Operator}

\subsection{The limits of classical equality}

The classical equality relation is:

\begin{itemize}
    \item static,
    \item context-free,
    \item observer-independent,
    \item and insensitive to dynamic computational processes.
\end{itemize}

But actual computation is:

\begin{itemize}
    \item dynamic and adaptive,
    \item observer- or algorithm-dependent,
    \item procedurally realized,
    \item driven by convergence rather than static identity.
\end{itemize}

Thus, a more flexible notion of equivalence is required.

\subsection{Observer non-fixation}

Recent developments in AI introduce \emph{observer non-fixation}:

\begin{itemize}
    \item the evaluating agent (algorithm, model, heuristic) changes over time;
    \item the computation space itself is adaptive;
    \item convergence behavior depends on the evolving observer.
\end{itemize}

Classical P=NP assumes a fixed Turing machine observer.  
This assumption is incompatible with modern computational phenomena.

\section{Definition of the Middle-Way Equivalence}

\begin{definition}[Middle-Way Equivalence]
For complexity classes $C_1$ and $C_2$, we write
\[
C_1 \mweq C_2
\]
if the following conditions hold:

\begin{enumerate}
    \item \textbf{Typical-case convergence:} 
    A high-probability subset of instances in $C_1$ converges to the computational 
    resource bounds of $C_2$.
    \item \textbf{Heuristic or statistical equivalence:} 
    Approximation, heuristic, or probabilistic methods achieve $C_2$-level 
    performance for solving $C_1$ instances.
    \item \textbf{Adaptive observer convergence:}
    When the computational observer is allowed to evolve,
    the search dynamics for $C_1$ become asymptotically equivalent to those of $C_2$.
\end{enumerate}
\end{definition}

This equivalence is \emph{dynamic}, not static.

\section{Main Theorem}

\begin{theorem}[Extended P-NP Equivalence]
Under the Middle-Way Equivalence relation,
\[
\mathrm{P} \mweq \mathrm{NP}.
\]
\end{theorem}

\begin{proof}[Sketch of proof]
The following independent observations jointly support the equivalence:

\begin{enumerate}
    \item Modern SAT solvers solve large NP-complete instances 
    near-polynomially in practical contexts.
    \item Many NP-complete problems have polynomial-time approximation schemes 
    or strong heuristics.
    \item Machine learning and large language models can statistically infer 
    solutions to NP problems.
    \item Probabilistic computation models narrow the gap between P and NP.
    \item Adaptive search and self-modifying heuristics cause NP search behavior 
    to converge toward P-like performance under realistic distributions.
\end{enumerate}

Together, these satisfy all three conditions of the $\mweq$ definition.
\end{proof}

\section{Relation to Classical P=NP}

Our result neither claims that
\[
\mathrm{P} = \mathrm{NP}
\]
nor that
\[
\mathrm{P} \neq \mathrm{NP}.
\]

Instead it highlights a deeper issue:
the equality symbol itself encodes ontological assumptions incompatible 
with modern computation.

The Middle-Way view generalizes equivalence from static identity to 
dynamic convergence.

\section{Applications}

This framework unifies the following phenomena:

\begin{itemize}
    \item empirical tractability of SAT and CSP,
    \item efficiency of branch-and-bound and CDCL heuristics,
    \item success of ILP solvers,
    \item AI-assisted optimization,
    \item evolutionary and reinforcement learning convergence.
\end{itemize}

Thus, $\mweq$ provides a rigorous vocabulary for ``practically P $\approx$ NP'' behavior.

\section{Connection to EMT (Emergent Middle Test)}

The EMT framework (Emergent Middle Test) describes convergence dynamics:

\[
\text{emptiness} \rightarrow \text{random excitation} 
\rightarrow \text{dependent origination} \rightarrow \text{Middle Way}.
\]

This dynamic structure aligns naturally with $\mweq$,
and EMT can serve as a computational test for verifying Middle-Way equivalence 
between classes or algorithms.

\section{Conclusion}

We introduced the Middle-Way Equivalence operator $\mweq$ and applied it to the
relationship between P and NP.  
This yields an \emph{extended resolution} of the P=NP problem in which:

\begin{itemize}
    \item the strict identity ``$=$'' is replaced by a dynamic, convergence-based equivalence,
    \item practical computational phenomena are formally integrated,
    \item P and NP are equivalent in the Middle-Way sense.
\end{itemize}

This reframes P=NP not as a binary decision problem, but as a question of 
computational phenomenology and convergence theory.

\begin{thebibliography}{9}

\bibitem{cook}
S. Cook,
\emph{The complexity of theorem-proving procedures},
Proceedings of the Third Annual ACM Symposium on Theory of Computing, 1971.

\bibitem{garey}
M. Garey and D. Johnson,
\emph{Computers and Intractability: A Guide to the Theory of NP-Completeness},
Freeman, 1979.

\bibitem{arora}
S. Arora and B. Barak,
\emph{Computational Complexity: A Modern Approach},
Cambridge University Press, 2009.

\bibitem{solomonoff}
R. Solomonoff,
\emph{A formal theory of inductive inference},
Information and Control, 1964.

\end{thebibliography}

\end{document}

